
\documentclass[preprint,12pt,authoryear]{elsarticle}


%% For including figures, graphicx.sty has been loaded in
%% elsarticle.cls. If you prefer to use the old commands
%% please give \usepackage{epsfig}

%% The amssymb package provides various useful mathematical symbols
\usepackage{amssymb}
%% The amsmath package provides various useful equation environments.
\usepackage{amsmath}
\usepackage{float}
\usepackage{hyperref}
\usepackage{amsthm}
\usepackage{xcolor}
\usepackage{booktabs}
\usepackage{tablefootnote}
\usepackage{comment}
\usepackage[noend]{algpseudocode}
\usepackage[ruled,vlined]{algorithm2e}


\usepackage{graphicx}
\definecolor{black}{RGB}{0,0,139}

%% The amsthm package provides extended theorem environments
%% \usepackage{amsthm}

%% The lineno packages adds line numbers. Start line numbering with
%% \begin{linenumbers}, end it with \end{linenumbers}. Or switch it on
%% for the whole article with \linenumbers.
%% \usepackage{lineno}

\journal{Spatial and Spatio-temporal Epidemiology}

\begin{document}

\begin{frontmatter}


\title{BARTSIMP: flexible spatial covariate modeling and prediction using Bayesian additive regression trees} %% Article title



%% Abstract
\begin{abstract}
%% Text of abstract
Prediction is a classic challenge in spatial statistics and the inclusion of spatial covariates can greatly improve predictive performance when incorporated into a model with latent spatial effects. It is desirable to develop flexible regression models that allow for nonlinearities and interactions in the covariate specification. Existing machine learning approaches that allow for spatial dependence in the residuals fail to provide reliable uncertainty estimates. In this paper, we investigate the combination of a Gaussian process spatial model with a Bayesian Additive Regression Tree (BART) model. The computational burden of the approach is reduced by combining Markov chain Monte Carlo (MCMC) with the Integrated Nested Laplace Approximation (INLA) technique. We study the performance of the method first via simulation. We then use the model to predict anthropometric responses in Kenya, with the data collected via a complex sampling design. In particular, household survey data are collected via stratified two-stage unequal probability cluster sampling, which requires special care when modeled. 
\end{abstract}

%%Graphical abstract
%\begin{graphicalabstract}
%\includegraphics{grabs}
%\end{graphicalabstract}

%%Research highlights
%\begin{highlights}
%\item Research highlight 1
%\item Research highlight 2
%\end{highlights}

%% Keywords
\begin{keyword}
%% keywords here, in the form: keyword \sep keyword
{Bayesian additive regression trees} \sep {Covariate modeling} \sep {Spatial prediction} \sep {Integrated nested Laplace approximation} \sep {Survey sampling}
%% PACS codes here, in the form: \PACS code \sep code

%% MSC codes here, in the form: \MSC code \sep code
%% or \MSC[2008] code \sep code (2000 is the default)

\end{keyword}

\end{frontmatter}

\section{Introduction}

There has been an increased interest in using covariates for spatial data modeling \citep{uwiringiyimana2022bayesian, macharia2019sub}, with previous studies showing that the inclusion of influential spatial covariates can lead to improved prediction accuracy \citep{lindstrom2014flexible, ZERAATPISHEH2022105723}. It is straightforward to include covariates in a linear model, within the linear predictor of spatial random effects models, but the extension to more flexible covariate models is more difficult. Flexible modeling is desirable to leverage nonlinearities and interactions which can lead to improved predictive performance.


Currently, Gaussian random field (GRF) models are commonly used as a modeling framework for capturing spatial correlations, with wide application, for example, in population health modeling \citep{diggle2019model}. {For data that arise from complex survey designs, fully Bayesian modeling approaches have been explored, see for example \cite{chan2020bayesian, banerjee2024finite} and \cite{finley2024models}.} A common approach to computation in spatial modeling is Integrated nested Laplace approximation (INLA) \citep{rue2009approximate}, particularly in a low-and-middle-income countries (LMIC) context \citep{utazi2018high, burstein2019mapping}. INLA is a powerful tool for carrying out Bayesian inference, but requires the predictors to have a linear form, which cannot capture multivariate nonlinearities, including  interactions, which may lead to reduced predictive performance. We do note that it is simple to include spline terms, since such models can be expressed in linear form, and also random walk models, both of which can simply capture univariate nonlinearities.

{There are many machine learning methods that allow for flexible covariate modeling, many of which have been applied to describing 
variation in health and demographic indicators , see for example \cite{bosco2017exploring} and \cite{burke2021using}. Particular methods that have been utilized include the convolutional neural network \citep{jean2016combining}, classification trees \citep{bilton2017classification}, stacked generalization \citep{davies2016optimal, osgood2018mapping} and random forests \citep{ren2018}. However, these approaches suffer from drawbacks. {A number of machine learning approaches in spatial modeling ignore the spatial dependence \citep{georganos2021geographical} while others (including the aforementioned references) either avoid giving uncertainty intervals or do not correctly propagate uncertainty, as is required for valid inference \citep{daw2023reds,wakefield2019estimating}}}. 

Bayesian Additive Regression Trees (BART) \citep{chipman2010bart} provide reliable Bayesian inference by specifying a prior distribution on the `sum-of-trees' structure that flexibly models the covariates. Previous applications of BART on spatial data modeling problems include \cite{muller2007spatially}, who proposed a spatial BART model based on a conditional autoregressive (CAR) model, and \cite{Krueger2020}, who suggested using a matrix exponential spatial specification \citep{lesage2007matrix} as an extension of CAR.  \cite{tan2019bayesian} review extensions to the basic BART  model, including a description of the approach due to \cite{muller2007spatially}.  However, compared to the GRF that we use, these models specify a simple variance-covariance structures and are designed for area-level data. The model considered by \cite{spanbauer2021nonparametric} assumes a more generic covariance structure under a non-spatial setting, but is limited to models with few random effects and does not scale well to large spatial datasets. As a result, we note that the large numbers of random effects in continuous spatial models, along with the strong dependence in the posterior, lead to computational difficulties and pose serious challenges for incorporating BART into spatial models.  

In this paper, we propose \textbf{BARTSIMP}, which is shorthand for `\textbf{BART} for \textbf{S}patial \textbf{I}NLA \textbf{M}odeling and \textbf{P}rediction', as a GRF spatial Bayesian modeling framework with a flexible covariate regression model. Our model leverages the flexibility in BART to capture the nonlinearity and interactions across covariates, while also recognizing the complex spatial correlation structure in the residuals which is modeled by the GRF. To ease computation, we use the INLA-within-MCMC method \citep{gomez2018markov} to design a Metropolis-within-Gibbs sampler that integrates out the random effects using INLA and then performs MCMC updates on the remaining parameters. { This model is the first to simultaneously model the covariate effects through BART and the spatial effects through a GRF, while producing Bayesian uncertainty estimates (credible intervals). Our spatial BART model is also the first to use INLA as an approximate approach to reduce the computational burden.}

We organize the paper as follows. We will describe our motivating example, which is spatial prediction for child anthropometric data, in Section \ref{sec:mot}. The model formulation is in Section \ref{sec:md} and we describe the Metropolis-within-Gibbs sampler which we use for model implementation in Section \ref{sec:comp}. We apply our model to simulated data experiments in Section \ref{sec:sim} and return to the Kenya data in Section \ref{sec:app}. Section \ref{sec:conc} concludes the paper with a discussion.

\section{Motivating Dataset}
\label{sec:mot}

Our study is motivated by child undernutrition data collected in the 2014 Kenya Demographic and Health Survey (DHS). Specifically, we are interested in wasting for children under the age of five, which is measured using the weight-for-height Z-score (WHZ) metric. The Z-score can be interpreted as the number of units that the weight for a child is higher or lower than the median value, compared to all children of the same height in the population \citep{mei2007standard}. For example, wasting is defined as a WHZ score below $-2$ \citep{kassie2019exploring}. 

The data is collected via stratified two-stage cluster sampling. In the first stage, 1584 enumeration areas (EAs) were sampled across 92 strata. The 92 sampling strata are defined according to the 47 counties and the urban-rural status of the county (with Nairobi and Mombasa counties only having urban areas). At the second stage, 25 households were selected from each of the EAs, to give 40,300 households in total and yielding 20,977 child WHZ measurements. 

We use geospatial covariates on the raster level, with a list of covariate descriptions in Table \ref{tab:cov} {(we have also included an extended version of Table \ref{tab:cov} including the range, types of transformation and range after transformation in Table 4 of the Supplementary Materials)}. In our example, for simplicity, we did not include the urban-rural status as a covariate in our model, which may lead to some bias due to oversampling of urban clusters but we believe that the inclusion of population density provides some protection.

\begin{table}[H]
\small
\centering
\caption{DHS Geospatial covariates used in our model.\\}
\begin{tabular}{lll}
\hline
\textbf{Variable} &
  \textbf{Years} &
  \textbf{Source} \\ \hline
Population Density &
  2014 &
  \begin{tabular}[c]{@{}c@{}}WorldPop\tablefootnote{http://www.worldpop.org.uk/data/get\_data/.} \end{tabular} \\ \hline
Night Time Light &
  2013 &
  \begin{tabular}[l]{@{}l@{}}National Centers for\\ Environmental Information\tablefootnote{https://ngdc.noaa.gov/eog/dmsp/downloadV4composites.html.} \end{tabular} \\ \hline
Vegetation Index &
  2013-2014 &
  \begin{tabular}[l]{@{}l@{}}NASA EOSDIS Land Processes\\ DAAC\tablefootnote{https://lpdaac.usgs.gov/dataset\_discovery/modis/modis\_products\_table/mod13a3\_v006.}\end{tabular} \\ 
\hline
Average Temperature &
  \begin{tabular}[l]{@{}l@{}}Mean value over\\ the period\\ 1970--2000\end{tabular} &
  \begin{tabular}[c]{@{}c@{}}WorldClim\tablefootnote{http://worldclim.org/version2.}\end{tabular} \\ \hline 
Precipitation &
  \begin{tabular}[l]{@{}l@{}}Mean value over\\ the period\\ 1970--2000\end{tabular} &
  \begin{tabular}[c]{@{}l@{}}WorldClim\tablefootnote{http://worldclim.org/version2.}\end{tabular} \\ \hline
Access to Nearest City &
  2000 &
  \begin{tabular}[l]{@{}l@{}}Joint Research Centre of the\\ European Commission\tablefootnote{http://forobs.jrc.ec.europa.eu/products/gam/download.php.}\end{tabular} \\ \hline
\end{tabular}
\label{tab:cov}
\end{table}

The aim of the analysis is to provide predictions of WHZ at various administrative levels. To introduce some terminology, Admin 0 denotes the national level, Admin 1 one below that (for example, states in the United States) and Admin 2 one below that (for example, counties in the United States). In Kenya, the Admin 1 level contains 47 counties, while Admin 2 contains 290 constituencies. Hence, our objective is a problem in small area estimation (SAE). If there are sufficient data in each area, a weighted  estimator can be used. Such estimates account for the design, but estimates in each area are based on data from that area only (and in the SAE literature are referred to as direct estimates) But often there are insufficient area-based data and models must be introduced. In an area-based model \citep{fay1979estimates} the weighted estimate is modeled and random effects are introduced. In a unit-level model, individual level data are modeled \citep{battese1988error}, and this is the path we follow. \cite{rao2015small} provides a comprehensive overview of SAE. For the Kenya data, Figure \ref{fig:direct} gives Admin 1, weighted estimates along with the widths of 90\% (asymptotic) confidence intervals (based on a variance estimate that accounts for the complex design). In the left panel, we see  higher levels of WHZ in the west and south, with lower, less desirable, outcomes in the north and east. In the right panel, we see that the accompanying uncertainty measures are wide, so that the weighted estimates are producing imprecise estimates{, especially within the areas in the northwest region, making it difficult to interpret the weighted estimates.} We aim to reduce uncertainty using spatial smoothing and covariate modeling. We emphasize that we model at the point level, but the final deliverable for policy making is at the area-level, so that we are required to aggregate from points to areas.


\begin{figure}[H]
\centering
\includegraphics[width = \linewidth]{direct_admin1_slice_1_aggregated.pdf}
\caption{Maps of Admin 1 level WHZ weighted estimates (left) and 90\% confidence interval widths (right).}
\label{fig:direct}
\end{figure}

\section{Model Description}
\label{sec:md}

To describe the data framework and introduce notation, we will outline a conventional traditional spatial model with a linear model in Section \ref{sec:LGM}. We will then introduce BART in Section \ref{sec:BART} and propose our spatial version in Section \ref{sec:SBART}. We describe the priors in Section \ref{sec:SBARTprior}.

\subsection{Model formulation and preliminaries}
\label{sec:LGM}

We consider the spatial dataset $(\mathbf{Y}, \boldsymbol{x})$, where $\mathbf{Y}$ is the collection of observed outcomes and $\boldsymbol{x}$ is the covariate matrix. Let $n$ be the number of spatial locations observed (in our example this corresponds to clusters), noting that we will allow multiple outcomes (these are children in our case) to be observed at the same location so that $\boldsymbol{y} = \left(\boldsymbol{y}_1, ..., \boldsymbol{y}_n\right)^\intercal$ is a jagged array of vectors, with $n_i$ being the number of observations in cluster $i$ and $\boldsymbol{y}_i = \left(y_{i1}, \dots, y_{in_i}\right)^\intercal$ being the vector of all observations at the same location. Specifically, we let $y_{ik}:=y_{k}(\boldsymbol{s}_i)$ be the $k$-th response observed at (two-dimensional) location $\boldsymbol{s}_i$, $i = 1, \dots, n$. We let $P$ be the number of covariates (including the intercept), so that $\boldsymbol{x}_i = (x_{i1}, \dots, x_{iP})$ is the vector of covariates observed at location $\boldsymbol{s}_i$ where $x_{ip}$ is the value of covariate $p$ observed at cluster $i$ for $p = 1, \dots, P$ (for simplicity of notation, we assume that observations at the same spatial unit share the same covariate values). We assume a Gaussian error model for the observations with measurement error  variance $\sigma_e^2$. The assumption of constant variance across locations is important, and we should endeavor to check this. We let the latent random spatial component at location $\boldsymbol{s}_i$ be denoted ${z}(\boldsymbol{s}_i)$, $i=1,\dots,n$, and assume that is follow a GRF with Mat\'ern covariance function over the region. 


A simple linear mixed effects model with a spatial random field can be written as: %We consider the following linear mixed effects model with spatial random fields:
{
\begin{align} 
y_{ik} \mid \boldsymbol{x}_{i}, \boldsymbol{\beta},  \boldsymbol{z}, \sigma_e^2 &\stackrel{i.i.d}{\sim} N\left(\boldsymbol{x}_i^{\top}\boldsymbol{\beta} +z_i, \sigma_e^2\right), \label{eq:model0} \\ 
\boldsymbol{z} \mid \boldsymbol{\psi}  &\sim G F(0, \sigma_m^2  \boldsymbol{\Sigma}),
\label{eq:model}
\end{align}
$k=1,\dots,n_i$, $i=1,\dots,n$.
The  
$n \times n$ matrix $\mathbf{\Sigma}$ is a Mat\'ern correlation matrix and $\sigma_m^2$ is the spatial variance. The correlation matrix has elements,
\begin{equation} 	
\Sigma_{i j} =\operatorname{Corr}_M\left(z\left(\boldsymbol{s}_i\right), z\left(\boldsymbol{s}_j\right)\right)=\frac{2^{1-\upsilon} }{\Gamma(\upsilon)}\left(\kappa\left\|\boldsymbol{s}_i-\boldsymbol{s}_j\right\|\right)^\nu K_\nu\left(\kappa\left\|\boldsymbol{s}_i-\boldsymbol{s}_j\right\|\right),
\label{eq:matern}
\end{equation} 
for $i,j=1,\dots,n$, $i \neq j$, and 
with $\kappa$ being the scale parameter of the GRF and $K_\upsilon(\cdot)$ is the modified Bessel function. The smoothness parameter $\upsilon$ is usually fixed \textit{a priori}, because data sparsity makes it difficult to estimate. {For the rest of the paper, we set the value of $\upsilon$ to the default value used in \citep{lindgren2011explicit}.} There is a one-to-one relationship between the scale parameter $\kappa$ and the range parameter (which is the more commonly interpreted parameter) with $\rho=\sqrt{ 8 \upsilon}/\kappa$. We let $\boldsymbol{\psi} = (\sigma_m^2, \rho)$ be the set of Mat\'ern covariance parameters.}  The model is completed with priors on the regression parameters $\boldsymbol{\beta}$ and a prior on the variance components,
\begin{equation}\label{eq:varprior}
\pi( \sigma_e^2,\sigma_m^2,\kappa).
\end{equation}
Details on the particular prior we use are postponed until Section \ref{sec:SBARTprior}.

For latent Gaussian models (LGMs), a closed-form expression for $\pi(\boldsymbol{z}, \boldsymbol{\beta}, \sigma_e^2, \sigma_m^2, \kappa \mid \boldsymbol{y})$ is computationally expensive for large $n$, due to the computation of the inverse and determinant of the dense $n \times n$ covariance matrix. Efficient MCMC sampling algorithms are difficult to construct for spatial models because of the strong dependence between parameters \citep{knorr2002block}. \cite{rue2009approximate} showed that the INLA method, originally an approximation Bayesian inference procedure based on Laplacian approximations and numerical integration rules, is available for LGMs. The INLA approach provides approximations for the marginal posterior distributions of the latent Gaussian random field and the hyperparameters, along with a reliable estimate for the marginal distribution $\pi(\boldsymbol{y})$ \citep{hubin2016estimating}.

For models with continuously indexed GRFs, \cite{lindgren2011explicit} showed that a discrete Gaussian Markov random field (GMRF) can be used to approximate the continuously indexed data GRF generated through a Mat\'ern covariance function, via a stochastic partial differential equation (SPDE) approach. This not only reduces the computation complexity from $\mathcal{O}(n^{3})$ to $\mathcal{O}(n^{3/2})$, but also enables model inference to be carried out through INLA. 


\subsection{The BART model}
\label{sec:BART}
%{
%
%The notation of $M$ and $\boldsymbol{\mu}$: I originally used $M$ to represent the node values, but then realized that it would be hard to use it to express the `matrix form'. Some other literature also used $\boldsymbol{\mu}$ as an alternative notation and I wonder if we should switch to that instead; $\psi = (\sigma_m^2, \rho)$ is the vector of Mat\'ern covariance hyperparameters. 
%}

%\begin{itemize}
%	\item Explain the components of BART: the trees are a set of decision rules and the node values are the model fit for the terminal nodes 
%	\item Explain interpretations of the model parameters related to BART
%	\item A lot of papers have simple examples, maybe do one too
%\end{itemize}
{
The Bayesian Additive Regression Trees (BART) model \citep{chipman2010bart} can be viewed as the summation of a fixed number of highly flexible nonparametric regression functions, where each function {$g(\cdot ; T_j, \boldsymbol{\mu}_j): \mathbb{R}^p \rightarrow \mathbb{R}, ~ j = 1, \dots, m$} is a random function that maps the $P$-dimensional covariate space onto the real line, with a total number of $l = 1, \dots, m$ functions. The component functions are defined by two sets of parameters, the binary tree structure $\mathbf{T}=\left(T_1, \ldots, T_m\right)$ and the terminal node values $\boldsymbol{\mu} = (\boldsymbol{\mu}_1, \ldots, \boldsymbol{\mu}_m)$. Each binary tree can be viewed as a set of decision rules that splits the covariate space into a finite number of regions and at any internal node, for a given tree, a `splitting variable' is chosen from the $P$ covariates and a threshold value is chosen that splits the current region into two sub-regions. A fitted value is then assigned to the corresponding region at each end node, according to the terminal node values $\boldsymbol{\mu}$. By splitting according to different covariate variables and threshold values, BART is capable of modeling complex nonlinear relationships and interactions among covariates. Furthermore, a regularization prior that penalizes tree complexity in order to avoid model overfitting is assigned to the tree structure parameters so that the trees functions are `weak learners' \citep{chipman2010bart}. 

\subsection{Sampling Model and Latent Field}
\label{sec:SBART}
%\begin{itemize}
%	\item Combine the previous two sections together, state the sampling model we will be using
%\end{itemize}
The sampling model for BARTSIMP is a combination of the linear mixed effects model described in Section \ref{sec:LGM} with the BART terms of Section \ref{sec:BART} replacing the linear model. Substituting the linear covariate effects with the sum-of-trees model gives: 
%We consider the spatial dataset $(\mathbf{Y}, \mathbf{X})$, where $\mathbf{Y}$ is the collection of observed outcomes and $\mathbf{X}$ is the covariate matrix. Let $n$ be the number of total spatial units in the dataset(for example, clusters in a survey), as we allow multiple outcomes to be observed at the same spatial unit, $\mathbf{Y} = \left(\boldsymbol{y}_1, ..., \boldsymbol{y}_n\right)^\intercal$ is a jagged array of vectors, with $n_i$ being the number of observations in unit $i$ and $\boldsymbol{y}_i = \left(y_{i}^{1}, \dots, y_{i}^{n_i}\right)^\intercal$ being the vector of all observations in the same place. Let $P$ be the number of covariates and $x_{ip}$ be the value of covariate $p$ observed at unit $i$ (we assume that observations within the same spatial unit shares the same covariate values). [Let $\mathbf{T} = (T_1, ..., T_m)$ be the collection of $m$ trees, and $\boldsymbol{\mu} = (M_1, ..., M_m)$ be the set of terminal node values.] We consider the following BART model with spatial random fields:
\begin{align*}
	y_{ik} \mid z_i, \boldsymbol{x}_{i p}, \mathbf{T}, \boldsymbol{\mu}, \sigma_e^2 \stackrel{i . i . d}{\sim} N\left(\sum_{l=1}^m g\left(\boldsymbol{x}_{i p} ; T_l, \boldsymbol{\mu}_l\right)+z_i, \sigma_e^2\right)
\end{align*}
with the distribution of $\boldsymbol{z}$ and Mat\'ern covariance being as previously defined in Equations \eqref{eq:model} and \eqref{eq:matern}.

}

\subsection{The Priors}
\label{sec:SBARTprior}
We assume a priori independence between the tree parameters $(\mathbf{T}, \boldsymbol{\mu})$, the residual variance parameter $\sigma_e^2$ and the spatial hyperparameters $\boldsymbol{\psi}$:
\begin{align*}
	p\left(\mathbf{T}, \boldsymbol{\mu}, \sigma_e^2, \boldsymbol{\psi}\right)=\left[\prod_{j=1}^m p\left(\boldsymbol{\mu}_j \mid T_j\right) p\left(T_j\right)\right] p\left(\sigma_e^2\right) p\left (\boldsymbol{\psi}\right).
\end{align*}
{For the tree parameters and the residual variance parameter, we follow the prior specifications in \cite{chipman1998bayesian} and decompose the tree structure prior into three hierarchical parts: the probability of a node being non-terminal, the probability of choosing one of the covariates as the splitting variable for a non-terminal node and the probability of choosing a splitting value given the chosen splitting variable. For a tree node at depth $d$, the probability of it being a non-terminal node is given the form $\alpha(1+d)^{-\beta}$, with $\alpha \in(0,1), \beta \in[0, \infty)$. In the  examples in the paper, we defer to the default settings in \cite{chipman2010bart} and choose $\alpha = 0.95$ and $\beta = 2$. Further, we assume that the splitting variable is uniformly chosen among all covariates for each internal node, and the splitting value is uniformly chosen from all distinct values of the selected splitting variable.} We assume independent and identically distributed normal priors for the terminal node values $\boldsymbol{\mu}_1, \dots, \boldsymbol{\mu}_m$ given their corresponding tree structure $T_j$. The mean and variance for the normal prior are chosen such that each tree will only function as a `weak learner' that contributes a small part in the model; we refer to \cite{chipman2010bart} for further details. 

For the residual variance parameter $\sigma_e^2$, we take the default choice in \cite{chipman2010bart} and specify a scaled inverse-Gamma distribution $\sigma_e^2 \sim \nu \lambda / \chi_\nu^2$, where $\nu$ is the degrees-of-freedom. The values of the hyperparameters $(\nu, \lambda)$ are chosen through an exploratory data analysis procedure to give $\hat{\sigma}_r^2$ as an empirical guess of $\sigma_e^2$ through a tentative working model (for example, an SPDE model with all covariates as linear predictors), which matches the $q$-th quantile of the scaled inverse Gamma prior. Following \cite{chipman2010bart}, we let $(q, \nu) = (0.9, 3)$. Hence, the prior depends on the data in a weak fashion. 

Finally, we place a penalized complexity (PC) prior on the Mat\'ern hyperparameters $\boldsymbol{\psi}$, following \cite{Fuglstad19}. The joint PC prior for $\boldsymbol{\psi}$ is 
\begin{align*}
	p(\boldsymbol{\psi}) = \frac{d}{2}\tilde{\lambda}_1\tilde{\lambda}_2\rho^{-d/2 - 1} \exp \left(-\tilde{\lambda}_1 \rho^{-d/2} - \tilde{\lambda}_2 \sigma_m \right), \sigma_m >0, \rho > 0,
\end{align*}
where $d=2$, $\tilde{\lambda}_1 = -\log(\alpha_1)\rho_0^{d/2}$ and $\tilde{\lambda}_2 = -{\log (\alpha_2)}/{\sigma_0}$. The hyperparameters $(\alpha_1, \alpha_2, \rho_0, \sigma_0)$ guarantee that 
$\operatorname{P}(\rho < \rho_0) = \alpha_1$ and $\operatorname{P}(\sigma_m > \sigma_0) = \alpha_2$. We let $\alpha_1 = \alpha_2 = 0.5$ and set $(\rho_0, \sigma_m)$ to the crude estimates obtained from fitting a model with only an intercept and the spatial random field. Hence, we again specify a data dependent prior. 

{
\section{Computation}
\label{sec:comp}

\subsection{Overview on the Algorithm}
}
{For the standard BART model, an efficient Gibbs-type sampler is available under conjugate priors on the model parameters \citep{chipman1998bayesian}. However, the computation faces scalability challenges once the spatial random effect component is added into the model.} Under our model framework, a standard MCMC algorithm requires the sampling and storage of the spatial random field $\boldsymbol{z}$, which is computationally undesirable. Hence, our strategy is to integrate out $\boldsymbol{z}$ and operate on the corresponding marginal likelihood. However, the calculation of the marginal likelihood involves computing the determinant of large and dense variance-covariance matrices, which becomes prohibitive as the number of datapoints at distinct locations (clusters in the Kenya example) becomes large. As a result, we consider the `INLA-within-MCMC' technique proposed by \cite{gomez2018markov}. Specifically, the algorithm uses INLA to approximate the aforementioned marginal likelihood, which is then used to compute the Metropolis-Hastings acceptance ratio in the MCMC routine for the remaining model parameters. %uses splits the model parameters into two sets: the spatial component $\boldsymbol{z}$ and the non-spatial components $(\mathbf{T}, \boldsymbol{\mu}, \sigma_e^2, \psi)$, and INLA is used within the Metropolis-Hastings algorithm to generate posterior samples of the non-spatial components.
%In this section, we outline the algorithmic details of our approach based on the `INLA-within-MCMC' technique \citep{gomez2018markov}. Specifically, the algorithm splits the model parameters into two sets: the spatial component $\boldsymbol{z}$ and the non-spatial components $(\mathbf{T}, \boldsymbol{\mu}, \sigma_e^2, \psi)$, and INLA is used within the Metropolis-Hastings algorithm to generate posterior samples of the non-spatial components.

{

{We conclude this section with a brief overview of the INLA-within-MCMC method, and refer the readers to Section 7 in the Supplementary Materials for full details of the computation algorithm for BARTSIMP.} We let $\boldsymbol{\theta}$ denote the whole ensemble of hyperparameters and latent effects. Further, we consider the following partition $\boldsymbol{\theta} = \left(\boldsymbol{\theta}_{c}, \boldsymbol{\theta}_{-c}\right)$, where $\boldsymbol{\theta}_{-c}$ represents the set of parameters we would like to integrate out, and $\boldsymbol{\theta}_{c}$ being the set of parameters for which we will use a Metropolis-Hastings algorithm. 
The posterior distribution can be re-expressed as 
$$
\pi(\boldsymbol{\theta} \mid \boldsymbol{y}) \propto \pi\left(\boldsymbol{y} \mid \boldsymbol{\theta}_{-c}, \boldsymbol{\theta}_{c}\right) \pi\left(\boldsymbol{\theta}_{-c} \mid \boldsymbol{\theta}_{c}\right) \pi\left(\boldsymbol{\theta}_{c}\right) .
$$
Integrating $\boldsymbol{\theta}_{-c}$ from both sides gives,
\begin{align}
\label{eq:int_pos}
  \pi\left(\boldsymbol{\theta}_{c} \mid y\right) \propto \pi\left(\boldsymbol{y} \mid \boldsymbol{\theta}_{c}\right) \pi\left(\boldsymbol{\theta}_{c}\right).  
\end{align}
In the context of the Metropolis-Hastings algorithm, let $\boldsymbol{\theta}_c^{\star}$ be the potential value for the new iteration proposed by the transition kernel $q\left(\cdot\mid \boldsymbol{\theta_c}\right)$. A new $\boldsymbol{\theta}_c$ is proposed within each iteration of the MCMC routine. Plugging \eqref{eq:int_pos} into the expression of the acceptance probability, we have, % becomes,
\begin{align}
\label{eq:acceptance_ratio}
\alpha=\min \left\{1, \frac{\pi\left(\boldsymbol{y} \mid \boldsymbol{\theta}_{c}^{*}\right) \pi\left(\boldsymbol{\theta}_{c}^{*}\right) q\left(\boldsymbol{\theta}_{c} \mid \boldsymbol{\theta}_{c}^{*}\right)}{\pi\left(\boldsymbol{y} \mid \boldsymbol{\theta}_{c}\right) \pi\left(\boldsymbol{\theta}_{c}\right) q\left(\boldsymbol{\theta}_{c}^{*} \mid \boldsymbol{\theta}_{c}\right)}\right\}.
\end{align}
With the prior distribution $\pi(\cdot)$ 
and proposal distribution $q(\cdot \mid \cdot)$ for $\boldsymbol{\theta}_c$ specified, the computation bottleneck in \eqref{eq:acceptance_ratio} is to calculate $\pi\left(\boldsymbol{y} \mid \boldsymbol{\theta}_{c}\right)$ and $\pi\left(\boldsymbol{y} \mid \boldsymbol{\theta}_{c}^{\star}\right)$. In practice, it can be computationally expensive to obtain an exact likelihood. The INLA-within-MCMC approach provides an efficient method to provide approximations to $\pi\left(\boldsymbol{y} \mid \boldsymbol{\theta}_{c}\right)$ and $\pi\left(\boldsymbol{y} \mid \boldsymbol{\theta}_{c}^{\star}\right)$, by fitting models with R-INLA, with the values of $\boldsymbol{\theta_c}$ (or $\boldsymbol{\theta_c}^{\star}$) fixed. 
%In our model where $\boldsymbol{\theta} = (T_1, ..., T_m, M_1, ..., M_m, \sigma_e, \sigma_m, \rho, \mathbf{z})$, we designate that \\ $\boldsymbol{\theta}_c = (T_1, ..., T_m, M_1, ..., M_m, \sigma_e, \sigma_m, \rho)$ and $\boldsymbol{\theta}_{-c} = \mathbf{z}$. Essentially, we are integrating out all the latent spatial variables that grows with the size of the dataset. In the following sections, the residual variance parameter $\sigma_e^2$ will be replaced by $\sigma^2$ for brevity. Finally we denote $\psi$ as the set of SPDE hyperparameters $\psi = (\sigma_m^2, \kappa)$.


\subsection{R and C++ Integration}

We implemented the backbone of the MCMC algorithm for BARTSIMP based on the C++ code in the \texttt{BART} package \citep{spanbauer2021nonparametric}. In order to use the functions in the \texttt{R-INLA} package \citep{martins2013bayesian} to carry out the INLA computation during the Metropolis-Hastings adjustment step, we used Rcpp \citep{eddelbuettel2011rcpp, eddelbuettel2013seamless, eddelbuettel2018extending} as an interface between C++ and R. To provide an open-source computing software for BARTSIMP, we developed the R package \texttt{BARTSIMP}, with source code available at \url{https://github.com/AlexJiang1125/BARTSIMP}. 

\section{Simulation Experiments}
\label{sec:sim}

\subsection{Simulation setting}

In this section, we study several simulation scenarios in which the spatial and covariate signals have different strengths. We consider a $50 \times 50$ grid surface over a study region $[0,1] \times [0,1]$ and denote the set of grid cells as $G$. For the covariates associated with each grid cell, we independently generate two covariates from a uniform distribution on $[0,1]$. Let $x_{gp}, p = 1,2$, represents the $p$-th covariate for grid cell $g$ and let $\boldsymbol{x}_g=(x_{g1},x_{g2}), g \in G$. The true deterministic field evaluated at cell $g$, denoted $f(\boldsymbol{x}_g)$, is defined as:%a weighted average of the two fields mentioned above:
\begin{align*}
	f(\boldsymbol{x}_g) = (1-\omega) \boldsymbol{z}_g^{\star} + \omega f_{0}(\boldsymbol{x}_g), \qquad g \in G,
\end{align*}
{
where the `baseline' spatial field $\boldsymbol{z}^{\star}$ is generated from a GRF with Mat\'ern parameters $\kappa = 2.5$ (note that this corresponds to $\rho = \frac{\sqrt{8\upsilon}}{\kappa} \approx 1.13$), $\sigma_m^2 = 0.5$. 
We let $f_{0}(\boldsymbol{x}_g)$ be the `baseline' covariate surface. Depending on the structure in the covariate surface, we  have three different scenarios for $f_{0}(\boldsymbol{x}_g)$:
\begin{itemize}
    \item \textbf{Tree-structured}: $f_{0}(\boldsymbol{x}_g)$ is generated based on the following step function: 
    \begin{align*}
    f_{0}(\boldsymbol{x}_g) = \begin{cases}
    3 & x_{g1} < 0.5 \\
    0 & x_{g1}\geq 0.5, x_{g2} \geq 0.5 \\
    -2 & x_{g1}\geq 0.5, x_{g2} < 0.5
    \end{cases}.
    \end{align*}
    \item \textbf{Linear}: $f_{0}(\boldsymbol{x}_g)$ is a linear combination of $x_{g1}$ and $x_{g2}$:
    \begin{align*}
        f_{0}(\boldsymbol{x}_g) =2x_{g1} - x_{g2}.
    \end{align*}
    \item \textbf{Smooth}:  $f_{0}(\boldsymbol{x}_g)$ is a smooth function of $x_{g1}$ and $x_{g2}$:
    \begin{align*}
        f_{0}(\boldsymbol{x}_g) =\sin(2\pi x_{g1}) + \cos(2 \pi x_{g2}) .
    \end{align*}
\end{itemize}
}
%generated based on the {function:%Friedman function \citep{friedman2001greedy} as follows: 
%\begin{align*}
%f_{0g}(\boldsymbol{x}) = \sin \left(\pi x_{g1} x_{g2}\right)+2\left(x_{g3}-0.5\right)^2+2 %x_{g4}+ x_{g5}, ~~ \forall g \in G.
%\end{align*} 
The scalar $\omega$ is fixed at different values and can be interpreted as the proportion of `covariate signal' among the overall signal. Here we consider five different scenarios for $\omega$: $1$ (covariate signal only), $0.8$ (strong covariate signal), $0.5$ (medium covariate signal), $0.2$ (weak covariate signal) and $0$ (spatial signal only), which we denote as scenarios 1, 2, 3, 4 and 5. We then randomly select 250 cells from the grid and randomly sample one location uniformly within each cell, ending up with 250 spatial points over the study region, mimicking the spatial locations for the clusters. The number of observations for each spatial location is sampled from a uniform distribution over $\{5,6,7,8,9,10\}$. For each observation at a given spatial location, its value is defined as the sum of the deterministic field value (spatial field plus covariate signal)  from the grid cell it belongs to, and Gaussian random noise with $\sigma_e^2 = 1$. We simulate $10$ datasets for each scenario. {Note that the covariate and spatial surface are independently generated for each dataset, and not held constant. For the BARTSIMP and BART model, we used {ensembles of 20 trees, and collected 2,000 posterior samples after a burn-in period of 2,000.}  We set the model hyperparameters for the PC prior choices $(\alpha_1, \alpha_2, \rho_0, \sigma_m)$ and scaled inverse-Gamma prior choices $(\nu, q)$ to the default values described in Section \ref{sec:SBARTprior}. We also conducted a sensitivity analyses in Supplementary Materials Section 1 to show that the model estimation results are robust to the choice of hyperparameter values.

We compare BARTSIMP against several alternatives, which can be grouped into two categories depending on whether a BART component is included in the model. For the BART-related methods, we consider a standard BART model without a spatial latent field \textbf{(BART)} and a modified version of BARTSIMP, in which the likelihood function is calculated using the exact formula, instead of the SPDE approximation \textbf{(BARTSIMP-EXACT)}. We include this approach to assess whether the INLA approximation to the marginal distribution is accurate. For the non-BART methods, we consider a GMRF spatial model with Mat\'ern covariance function fitted by the SPDE approach, including all covariates in a linear model \textbf{(SPDE)} and a spatial SPDE model, with intercept only, i.e.,~no covariates, \textbf{(SPDE0)}. Additionally, we implemented an alternative version of the SPDE model, with a first-order random walk processes on the covariates \textbf{(SPDE-RW1)}.}  %Ideally we would simulate a greater number of datasets but even with shortcuts, the computation time is long. 




\subsection{Performance Criteria}
\label{sec:perfc}


{
 We examine the models via performance measures that evaluate both point and interval estimates, over the gridded surface:} %We focus on accuracy of interval estimates since this is clearly key for prediction and machine learning methods are often poor with respect to this aspect. 
\begin{itemize}
        {\item \textbf{Root mean squared error (RMSE)}: Let $\hat{f}^{(j)}(\boldsymbol{x}_g)$ be the model prediction for $f(\boldsymbol{x}_g)$ given by method $j$, the RMSE is defined as:
        \begin{align*}
            \text{RMSE}^{(j)} = \frac{1}{|G|}\sum_{g \in G}\left({\hat f^{(j)}(\boldsymbol{x}_g}) - f(\boldsymbol{x}_g)\right)^2, \qquad j =1,\ldots, 5.
        \end{align*}
        \item \textbf{Average interval length (AIL)}: Let $(L_{g,\alpha}^{(j)}, U_{g,\alpha}^{(j)})$ denote the $100 \times (1 - \alpha) \%$ prediction interval for $f(\boldsymbol{x}_g)$, available from method $j$. The AIL is defined as the average length of prediction intervals over the gridded surface:
        \begin{align*}
            \text{AIL}^{(j)} = \frac{1}{|G|}\sum_{g \in G}|U_{g, \alpha}^{(j)} - L_{g, \alpha}^{(j)}|, \qquad j =1,\ldots, 5
        \end{align*}}
	\item \textbf{Average coverage rate (ACR)}: 
    %Let $(L_{g,\alpha}^{(j)}, U_{g,\alpha}^{(j)})$ denote the $100 \times (1 - \alpha) \%$ prediction interval for $f(g)$ given by method $j$, 
    The ACR is defined as the coverage of $f(\boldsymbol{x}_g)$, for the interval, $(L_{g,\alpha}^{(j)}, U_{g,\alpha}^{(j)})$:
	\begin{align*}
		\text{ACR}^{(j)} = \frac{1}{|G|} \sum_{g \in G} \operatorname{I}\left(f(\boldsymbol{x}_g) \in L_{g,\alpha}^{(j)}, U_{g,\alpha}^{(j)})\right), \qquad j =1,\ldots, 5.
	\end{align*}
	For AIL and ACR we take $\alpha = 0.05$.
	\item \textbf{Average Interval score (AIS)}: The AIS is an integrated metric for prediction interval accuracy defined as: 
	\begin{multline*}
		\text{AIS}^{(j)} = \frac{1}{|G|} \sum_{g \in G} \left[ (U_{g,\alpha}^{(j)} -  L_{g,\alpha}^{(j)}) + \frac{2}{\alpha} \left(L_{g,\alpha}^{(j)} - f(\boldsymbol{x}_g) \right) \cdot \mathbf{1}(f(\boldsymbol{x}_g) < L_{g,\alpha}^{(j)})\right. + \\ \left.
		\frac{2}{\alpha} \left(f(\boldsymbol{x}_g) - U_{g,\alpha}^{(j)}\right) \cdot \mathbf{1}(f(\boldsymbol{x}_g) > U_{g,\alpha}^{(j)}) \right],  \qquad j =1,\ldots, 5.
	\end{multline*}
The AIS can be broken down into three parts, with the first part penalizing  wider intervals and the second and third parts penalizing low coverage rates \citealp{gneiting2007strictly}.

Among the measures, RMSE assesses the point prediction accuracy, while AIL, ACR and AIS focus on the uncertainty prediction accuracy. 
	
\end{itemize} 



\subsection{Results}

{Figures \ref{fig:modcomp} shows comparisons of RMSE, AIL, ACR and AIS across all six methods under the five scenarios, with the covariate surface generated from a tree-structured model.} First, we observe that BARTSIMP and BARTSIMP-EXACT perform similarly for all four metrics { under all scenarios and mechanisms}, suggesting that the approximation error from SPDE is negligible in model fitting. For point prediction accuracy, BART (the covariate-only model) has the lowest RMSE when there is only covariate signal, and both SPDE and SPDE0 (the spatial models) have the lowest RMSE when there is only spatial signal. Meanwhile, BARTSIMP (and BARTSIMP-EXACT) perform the best across all methods for the remaining scenarios where the true data is generated from a mixture of covariate and spatial signals. For uncertainty estimation, BARTSIMP performs best out of all methods in all scenarios, having coverage rates closest to the nominal coverage. As spatial signals become stronger, SPDE and SPDE0 tend to have lower interval widths and interval scores, and both BART and BARTSIMP tend to have wider interval widths and interval scores.

BARTSIMP and BARTSIMP-EXACT perform similarly, but their computation time differs. While the exact method is quicker for data with few observations, the approximation method is considerably more efficient as the number of observations grows. Attaining the nominal coverage is an issue for all approaches, particular for the stronger spatial signal cases, but we think this is just the reality of modeling point-level spatial data without large sample sizes. The same behavior of under-coverage was seen in the comprehensive examination of \cite{osgood2023statistical} (Figure 2 of this paper, in particular, clearly shows this phenomenon). 

{ The results for cases where the covariate surfaces are generated from a linear (or smooth) model are shown in Figure 13 and 14 in Supplementary Materials, Section 8. We note that the BART, BARTSIMP and BARTSIMP-EXACT still perform in a very similar fashion  to as they did in the  tree-structured case, where better prediction performances are associated with stronger covariate signals. For the case where the covariate surface is generated from a linear model, SPDE and SPDE-RW consistently outperform other methods in terms of most prediction metrics. This is not surprising, as the linear covariate effect can be correctly specified by these methods (however, note that all three BART-related methods have very similar RMSE when only the covariate signal is present in the data). As the model does not include covariate effects, SPDE0 performs worse than other methods when the covariate signal is strong, but has better performance as the spatial signal increases. 

When the covariate surface is generated from a smooth model, only SPDE-RW outperforms the other methods, as the smooth function cannot be correctly specified by models including SPDE. Note that BARTSIMP (and BARTSIMP-EXACT) still attain similar prediction performance when the covariate signal is relatively strong. 

Finally, we also conducted convergence analyses on the posterior samples from the simulation studies under all three scenarios, with results shown in the Supplementary Materials, Section 9.
}


\begin{comment}
\begin{table}[!ht]
\caption{RMSE for four different methods under four scenarios (covariate only, strong, medium and weak covariate signals). The averaged values over 10 replications are shown in the cells, with standard deviations in brackets.}
\centering
\begin{tabular}{@{}ccccc@{}}
\toprule
 &
  \begin{tabular}[c]{@{}c@{}}scenario 1\\ $(\omega = 1)$\end{tabular} &
  \begin{tabular}[c]{@{}c@{}}scenario 2\\ $(\omega = 0.8)$\end{tabular} &
  \begin{tabular}[c]{@{}c@{}}scenario 3\\ $(\omega = 0.5)$\end{tabular} &
  \begin{tabular}[c]{@{}c@{}}scenario 4\\ $(\omega = 0.2)$\end{tabular} \\ \midrule
BARTSIMP &
  \begin{tabular}[c]{@{}c@{}}1.260\\(0.051)\end{tabular} &
  \begin{tabular}[c]{@{}c@{}}1.266\\(0.042)\end{tabular} &
  \begin{tabular}[c]{@{}c@{}}2.558\\(0.019)\end{tabular} &
  \begin{tabular}[c]{@{}c@{}}2.042\\(0.016)\end{tabular} \\
BART &
  \begin{tabular}[c]{@{}c@{}}\textbf{1.064}\\(0.143)\end{tabular} &
  \begin{tabular}[c]{@{}c@{}}\textbf{0.787}\\(0.101)\end{tabular} &
  \begin{tabular}[c]{@{}c@{}}0.491\\(0.064)\end{tabular} &
  \begin{tabular}[c]{@{}c@{}}0.498\\(0.111)\end{tabular}\\
SPDE &
  \begin{tabular}[c]{@{}c@{}}1.241\\(0.195)\end{tabular} &
  \begin{tabular}[c]{@{}c@{}}0.925\\(0.135)\end{tabular} &
  \begin{tabular}[c]{@{}c@{}}0.734\\(0.170)\end{tabular} &
  \begin{tabular}[c]{@{}c@{}}0.670\\(0.222)\end{tabular} \\
SPDE0 &
  \begin{tabular}[c]{@{}c@{}}1.236\\(0.175)\end{tabular}&
  \begin{tabular}[c]{@{}c@{}}0.873\\(0.106)\end{tabular} &
  \begin{tabular}[c]{@{}c@{}}\textbf{0.459}\\(0.027)\end{tabular} &
  \begin{tabular}[c]{@{}c@{}}\textbf{0.183}\\(0.007)\end{tabular} \\ \bottomrule
\end{tabular}
\label{sim-methcomp}
\end{table}
\end{comment}

\begin{figure}
    \centering
    \includegraphics[width=\linewidth]{sim_study/sim_compare_tree.pdf}
    \caption{Root mean squared error (RMSE), average coverage rate (ACR), average interval length (AIL) and average interval score (AIS) for five different methods under five scenarios (covariate only, strong, medium, weak covariate signals and spatial signals only), when the covariate surface has a \textbf{tree-based structure}. The dots show the average values over 10 replications, and the vertical lines illustrate the 95\% Wald type interval computed from the standard deviation over 10 replications. The horizontal dashed line in the upper right figure illustrates the 95\% nominal coverage. The dashed horizontal line represents the nominal coverage rate (0.95). The horizontal positions of the dots and lines are slightly jittered to avoid overlapping, as some methods have very similar metrics.}
    \label{fig:modcomp}
\end{figure}

\begin{comment}

\begin{table}[!ht]
\caption{Average coverage rates (ACR) for four different methods (nominal coverage is 90\%) under four scenarios (covariate only, strong, medium and weak covariate signals). The averaged values over 10 replications are given in the cells, with standard deviations in brackets.}
\label{sim-methcomp2}
\begin{tabular}{@{}ccccc@{}}
\toprule
 &
  \begin{tabular}[c]{@{}c@{}}Scenario 1\\ $(\omega = 1)$\end{tabular} &
  \begin{tabular}[c]{@{}c@{}}Scenario 2\\ $(\omega = 0.8)$\end{tabular} &
  \begin{tabular}[c]{@{}c@{}}Scenario 3\\ $(\omega = 0.5)$\end{tabular} &
  \begin{tabular}[c]{@{}c@{}}Scenario 4\\ $(\omega = 0.2)$\end{tabular} \\ \midrule
BARTSIMP &
  \begin{tabular}[c]{@{}c@{}}\textbf{82}\%\\(0.016)\end{tabular} &
  \begin{tabular}[c]{@{}c@{}}79\%\\(0.023)\end{tabular} &
  \begin{tabular}[c]{@{}c@{}}59\%\\(0.005)\end{tabular} &
  \begin{tabular}[c]{@{}c@{}}78\%\\(0.007)\end{tabular} \\
BART &
  \begin{tabular}[c]{@{}c@{}}32\%\\(0.069)\end{tabular} &
  \begin{tabular}[c]{@{}c@{}}33\%\\(0.064)\end{tabular} &
  \begin{tabular}[c]{@{}c@{}}37\%\\(0.096)\end{tabular} &
  \begin{tabular}[c]{@{}c@{}}35\%\\(0.061)\end{tabular} \\
SPDE &
  \begin{tabular}[c]{@{}c@{}}98\%\\(0.000)\end{tabular} &
  \begin{tabular}[c]{@{}c@{}}\textbf{99}\%\\(0.000)\end{tabular} &
  \begin{tabular}[c]{@{}c@{}}\textbf{100}\%\\(0.000)\end{tabular} &
  \begin{tabular}[c]{@{}c@{}}\textbf{100}\%\\(0.000)\end{tabular} \\
SPDE0 &
  \begin{tabular}[c]{@{}c@{}}51\%\\(0.020)\end{tabular} &
  \begin{tabular}[c]{@{}c@{}}54\%\\(0.014)\end{tabular} &
  \begin{tabular}[c]{@{}c@{}}57\%\\(0.027)\end{tabular} &
  \begin{tabular}[c]{@{}c@{}}63\%\\(0.037)\end{tabular} \\ \bottomrule
\end{tabular}
\end{table}

\begin{table}[t]
\caption{Average interval score (AIS) rates for four different methods under four scenarios (covariate only, strong, medium and weak covariate signals). The averaged values over 10 replications are given in the cells, with standard deviations in brackets. (Low scores are preferred.)}
\label{sim-methcomp3}
\begin{tabular}{@{}ccccc@{}}
\toprule
 &
  \begin{tabular}[c]{@{}c@{}}scenario 1\\ $(\omega = 1)$\end{tabular} &
  \begin{tabular}[c]{@{}c@{}}scenario 2\\ $(\omega = 0.8)$\end{tabular} &
  \begin{tabular}[c]{@{}c@{}}scenario 3\\ $(\omega = 0.5)$\end{tabular} &
  \begin{tabular}[c]{@{}c@{}}scenario 4\\ $(\omega = 0.2)$\end{tabular} \\ \midrule
BARTSIMP &
  \begin{tabular}[c]{@{}c@{}}\textbf{0.286}\\(0.014)\end{tabular} &
  \begin{tabular}[c]{@{}c@{}}\textbf{0.290}\\(0.019)\end{tabular} &
  \begin{tabular}[c]{@{}c@{}}0.860\\(0.011)\end{tabular} &
  \begin{tabular}[c]{@{}c@{}}0.433\\(0.007)\end{tabular}\\
BART &
  \begin{tabular}[c]{@{}c@{}}0.503\\(0.108)\end{tabular} &
  \begin{tabular}[c]{@{}c@{}}0.380\\(0.077)\end{tabular} &
  \begin{tabular}[c]{@{}c@{}}0.224\\(0.050)\end{tabular} &
  \begin{tabular}[c]{@{}c@{}}0.224\\(0.052)\end{tabular}\\
SPDE &
  \begin{tabular}[c]{@{}c@{}}0.372\\(0.012)\end{tabular} &
  \begin{tabular}[c]{@{}c@{}}0.356\\(0.005)\end{tabular} &
  \begin{tabular}[c]{@{}c@{}}0.348\\(0.006)\end{tabular} &
  \begin{tabular}[c]{@{}c@{}}0.344\\(0.006)\end{tabular} \\
SPDE0 &
  \begin{tabular}[c]{@{}c@{}}0.462\\(0.069)\end{tabular} &
  \begin{tabular}[c]{@{}c@{}}0.314\\(0.038)\end{tabular} &
  \begin{tabular}[c]{@{}c@{}}\textbf{0.152}\\(0.008)\end{tabular} &
  \begin{tabular}[c]{@{}c@{}}\textbf{0.053}\\(0.003)\end{tabular} \\ \bottomrule
\end{tabular}
\end{table}
  
\end{comment}
\section{Application}
\label{sec:app}


In this section, we investigate the prediction performance of BARTSIMP, as compared to BART, SPDE and SPDE0, evaluated on the WHZ measurements from the 2014 Kenya DHS. 

\subsection{Cross-Validation Exercise}

To compare the prediction performances across different approacahes, we apply cross-validation and split the 1584 clusters in the 2014 DHS survey into training and test datasets of cluster sizes 1267 and 317, respectively. To preserve the stratification structure of the design in both data sets, we consider stratified sampling of the clusters such that the strata proportions in the training set roughly matches that in the test set. We repeated the procedure 10 times to reduce sampling variation caused by using a single data split. {We again use RMSE, ACR, AIL and AIS as performance criteria for all four methods, similar to Section \ref{sec:perfc}, and we use the same settings as before for the BART and BARTSIMP models. The algorithm took one day to run for BARTSIMP on Ubuntu 18.04, with 50GB of memory.

\begin{comment}
	We consider the following prediction performance measures:
\begin{itemize}
		\item \textbf{Bias}: To evaluate the deterministic estimation error of the methods, we can also evaluate the bias of the model estimate:
	\begin{align*}
		\text{Bias}^{j(k)} = \frac{1}{n} \sum_{i=1}^n \left(\hat{Y}_{i}^{j(k)} - y_i\right).
	\end{align*}
	\item \textbf{Average coverage rate (ACR):} Let $(L_{i,\alpha}^{j(k)}, U_{i,\alpha}^{j(k)})$ denote the $100 \times (1 - \alpha) \%$ prediction interval for $y_i^{(k)}$, the ACR is defined as the proportion of correct coverages of $y_i^{(k)}$:
	\begin{align*}
		\text{ACR}^{j(k)} = \frac{1}{n} \operatorname{I}\left(y_i^{(k)} \in L_{i,\alpha}^{j(k)}, U_{i,\alpha}^{j(k)})\right).
	\end{align*}
	\item \textbf{Average Interval score (AIS, \citealp{gneiting2007strictly})}: the AIS is an integrated metric for prediction interval accuracy is defined as follows: 
	\begin{multline*}
		\text{AIS}^{(j)} = \frac{1}{|G|} \sum_{g \in G} \left[ (U_{g,\alpha}^{(j)} -  L_{g,\alpha}^{(j)}) + \frac{2}{\alpha} \left(L_{g,\alpha}^{(j)} - f(g) \right) \cdot \mathbf{1}(f(g) < L_{g,\alpha}^{(j)})\right. + \\ \left.
		\frac{2}{\alpha} \left(f(g) - U_{g,\alpha}^{(j)}\right) \cdot \mathbf{1}(f(g) > U_{g,\alpha}^{(j)}) \right].
	\end{multline*}
	In our analysis, we let $\alpha = 0.10$.
\end{itemize} 
\end{comment}
Table \ref{fig:methcomp} shows ACR (nominal coverage is 95\%), RMSE, AIL and AIS for all approaches. The results are the average over 10 test datasets, with standard deviations in brackets. We see that BARTSIMP and BART dominate the other two methods in terms of having closest to nominal coverage.  The BARTSIMP model has the lowest AIS. The SPDE and SPDE0 models both have narrow intervals, but these produce low coverage, which results in poor AIS scores also.} 
\begin{comment}
Table \ref{fig:methcomp} shows ACR (nominal coverage is 95\%), RMSE, AIL and AIS for all approaches. The results are the averaged over 10 test datasets, with standard deviations in brackets. We see that BARTSIMP, BARTSIMP-EXACT and BART dominate the other two methods in terms of having closest to nominal coverage.  The BARTSIMP model has the lowest AIS.} 

 \begin{table}[ht]
\centering
\begin{tabular}{rlrrrrr}
  \hline
 & Model & ACR & RMSE & AIL & AIS \\ 
  \hline
1 & BARTSIMP & \textbf{98}\% (0.003) & 0.584 (0.024) & 6.882 (0.082) & \textbf{0.172} (0.002) \\
2 & BARTSIMP-EXACT & \textbf{97}\% (0.003) & 0.580 (0.026) & 6.680 (0.131) & 0.189 (0.004)\\ 
  3 & BART & \textbf{93}\% (0.005) & 0.590 (0.024) & 5.287 (0.028) & 0.182 (0.005)  \\ 
  4 & SPDE & 78\% (0.010) & \textbf{0.565} (0.023) & 3.136 (0.030) & 0.256 (0.009)\\ 
  5 & SPDE0 & 78\% (0.008) & 0.578 (0.025) & 3.156 (0.032) & 0.257 (0.010) \\ 
   \hline
\end{tabular}
\caption{Prediction performance measures over all five competing methods. The nominal coverage is 95\% and small values of AIS are preferred. The averages over $10$ test datasets are shown with standard deviations shown in brackets.}
\label{fig:methcomp}
\end{table}   
\end{comment}
\begin{table}[ht]
\centering
\caption{Prediction performance measures over all four competing methods. The nominal coverage is 95\% and small values of AIS are preferred. The averages over $10$ test datasets are shown with standard deviations shown in brackets.\\}
\begin{tabular}{rlrrrrr}
  \hline
 & Model & ACR & RMSE & AIL & AIS \\ 
  \hline
1 & BART-SIMP & \textbf{98}\% (0.003) & 0.584 (0.024) & 6.882 (0.082) & \textbf{0.172} (0.002) \\
  2 & BART & \textbf{93}\% (0.005) & 0.590 (0.024) & 5.287 (0.028) & 0.182 (0.005)  \\ 
  3 & SPDE & 78\% (0.010) & \textbf{0.565} (0.023) & \textbf{3.136} (0.030) & 0.256 (0.009)\\ 
  4 & SPDE0 & 78\% (0.008) & 0.578 (0.025) & 3.156 (0.032) & 0.257 (0.010) \\ 
   \hline
\end{tabular}
\label{fig:methcomp}
\end{table}
\subsection{Point and Areal Prediction}

In this section, we conduct spatial predictions on a grid surface over the study region and generate aggregated estimates at the Admin 1 and Admin 2 levels.
 The reason we produce estimates at the areal-level is that resource allocation and policy making decisions are made at the area-level.

At location $\boldsymbol{s}$, we let $\text{WHZ}(\boldsymbol{s})$ be the spatial surface of the height-for-weight Z-scores and $d_{5}(\boldsymbol{s})$ be the under-five population density. The areal level WHZ is a weighted average over the under-five population density, as the Z-scores were evaluated for children under age five. The  $d_5(\boldsymbol{s})$ values are obtained from WorldPop (\url{https://www.worldpop.org}). The WHZ for an administrative region $R_i$ is 
\begin{align}
\label{eq:sp}
	\text{WHZ}_{R_i} = \frac{\int_{R_i} \text{WHZ}(\boldsymbol{s}) d_5(\boldsymbol{s})}{\int_{R_i}  d_5(\boldsymbol{s})}, \qquad \qquad i = 1,2, \dots, m,
\end{align}
where $m$ is the number of administrative areas.
%where $\mathcal{R} = \{R_1, R_2, \dots\}$ is the set of administrative regions. 

We approximate the integrals in \eqref{eq:sp} by a weighted sum over observations on grid cells located at $\boldsymbol{s}_g$ over the regions, $g \in R_i$. Let $\text{WHZ}(\boldsymbol{s}_g)$ be the height-for-weight Z-score and $d_5(\boldsymbol{s}_g)$ be the under-five population density evaluated at grid cell $g$. The regional WHZ, approximated on the grid, is
\begin{align}
\label{eq:WHZ}
	{\text{WHZ}}_{R_i} \approx  \frac{\sum_{g \in R_i} {\text{WHZ}}(\boldsymbol{s}_g) d_5(\boldsymbol{s}_g) }{\sum_{g \in R_i} d_5(\boldsymbol{s}_g)}.
\end{align}
Using \eqref{eq:WHZ}, we can calculate the posterior mean and $95\%$ credible interval quantiles for regional ${\text{WHZ}}$ at all Admin 1 and Admin 2 areas, based on the four methods. As a comparison, we also consider the direct weighted areal-level estimates and 95\% (design-based) confidence intervals. Figure \ref{fig:admin1comps} shows the posterior median and 95\% credible/confidence intervals, derived from all five methods, and Figure \ref{fig:admin1areas} shows the predicted areal-level WHZ for all 47 Admin 1 areas in Kenya. The predicted posterior mean for WHZ given by BARTSIMP ranges from -1.10 to 0.05 over Admin 1 regions, showing that there is large within country variation in WHZ in Kenya. Among the 47 Admin 1 regions, Kiambu and Nairobi have the highest WHZ predictions -- these are the most populated counties in Kenya. Low WHZ scores occurred in areas such as Turkana, Mandera and Marsabit and these could be targeted for interventions. 

We see from Figures \ref{fig:admin1comps} and \ref{fig:admin1areas} that while all five methods give quantitatively similar overall patterns of the areal estimate across the regions, BARTSIMP and the spatial methods provide similar point estimates at the local level. BARTSIMP gives relatively wider credible interval lengths, which  is consistent with the simulations, in which we saw that  these intervals are more appropriate. Finally, BART does not yield reliable point estimates, and gives interval estimates that are far too narrow. {We also provide areal-level predictions on Admin 2 levels in Section 3 of the Supplementary Materials (results shown in Figures 4 to 11). Examination of these results show that the direct estimates for Admin 2 regions have more unreliable point estimates with much wider confidence intervals, due to insufficient samples observed in each region.}

{ Additionally, an interesting question is whether the GRF models can address the correlation within household. While it is theoretically straightforward to include an additional random household effect to account for dependence on observations in the same household, when we have attempted this in other projects with DHS data, there were identifiability issues, since it was hard to disentangle household and cluster effects.}


\begin{figure}[H]
\centering
%\includegraphics[width = \linewidth]{admin1_slice_1_aggregated_comparisons.pdf}
\includegraphics[width = .8\linewidth]{admin1_slice_1_aggregated_comparisons.pdf}
\caption{The Admin 1 level posterior median and 95\% credible/confidence intervals of WHZ for BART, BARTSIMP, SPDE0, SPDE and direct (weighted) estimates. The Admin 1 areas are arranged according to the predicted posterior mean given by BARTSIMP.}
\label{fig:admin1comps}
\end{figure}

\begin{figure}[H]
\centering
\includegraphics[width = .49\linewidth]{admin1_comp_med.pdf}
\includegraphics[width = .49\linewidth]{admin1_comp_ciw.pdf}
\caption{Maps of Admin 1 level WHZ posterior median (left) and 90\% credible interval lengths (right) for BARTSIMP, BART, SPDE and SPDE0.}
\label{fig:admin1areas}
\end{figure}





\subsection{Partial dependence}

In our example, it is also of  interest to study the marginal influence of each variable. Various measures of partial dependence have been proposed in the machine learning literature \citep{breiman2001random, goldstein2015peeking}.  For models based on BART, the partial dependence function \citep{friedman2001greedy} is a commonly used measure. Let $f(\cdot): \mathbb{R}^p \rightarrow \mathbb{R}$ be a multivariate function defined on $p$ variables $\boldsymbol{x}= (x_1, \dots, x_p) $. Let $\boldsymbol{x}_t$ denote the index of variables we wish to study, out of the $p$ variables, and let $\boldsymbol{x}_c = \boldsymbol{x} \slash {x}_t$ be their compliment. Suppose we have $n$ observations of such multivariate variables. The partial dependence function for $f(\cdot)$ with respect to ${x}_t$ is defined as 
\begin{align*}
	f^{pd}({x}_t) = \frac{1}{n} \sum_{i=1}^n f({x}_t, \boldsymbol{x}_{ic}),
\end{align*}
where $\boldsymbol{x}_{i.}, i = 1, \dots, n$ represents the $i$-th observation. For BARTSIMP and BART, we can calculate the partial dependence function of the sum-of-trees function for all six variables in our dataset. Figure \ref{fig:pd1} shows the partial dependence function and its $95\%$ credible interval based on {2000 MCMC samples for the BARTSIMP and BART methods. We observe that for all six variables, BARTSIMP and BART have similar partial dependence patterns that roughly resemble the pattern of the raw data. Substantively, holding all other variables constant, WHZ increases with increasing population density and precipitation, and decreases with increasing average temperature and access to nearest city. Note that BART has narrower credible intervals, especially when the covariate values are more extreme, which support our previous finding that BART tends to underestimate model uncertainty in this setting. We compare the results in Figure \ref{fig:pd1} to the posterior summary of the fixed effects, fitted from an SPDE model to the same dataset (see Table 5 in Supplementary Materials Section 6 for detailed estimation results). We note that for the SPDE model, the population density has a  positive association, while covariates including access to nearest city, average temperature have negative associations, this is in line with the results shown in Figure \ref{fig:pd1}. Additionally, we note that these results are plotted on the original scale of the covariates (e.g.,~before transformation), which caused the majority of the data points used to generate the line plots to be `clustered' towards the left end of the $x$-axis. An alternative partial dependence plot is shown in Supplementary Materials Section 5, where the data points are plotted on the transformed scale. }

%Note that BART has much narrower credible intervals, which support our previous finding that BART tends to underestimate model uncertainty in this setting. The association with VEGE is nonlinear but there is great uncertainty at low levels of access.

\begin{figure*}[!ht]
	\centering
	\includegraphics[width = \linewidth]{partial/res_o.pdf}
	\caption{Partial dependence function for all six covariates (population density, average temperature, night time light, vegetation index, precipitation and access to nearest city) and the 95\% pointwise credible interval provided by the BARTSIMP (blue) and BART (red) methods.}
	\label{fig:pd1}
\end{figure*}

\section{Discussion}
\label{sec:conc}

In this work, we have proposed BARTSIMP as a novel framework for flexible covariate modeling and prediction for spatial datasets. We incorporated the nonparametric nature of BART into continuous spatial models and leveraged the flexibility of BART to allow nonlinear covariate relationships and interactions. To implement the model we developed a sampling-based inference algorithm for BARTSIMP, based on the INLA-within-MCMC technique. 

BARTSIMP has a number of limitations which require more investigation. A cross-validation analysis showed that the BARTSIMP method yields average coverage rates closer to the nominal coverage compared to other methods while having poorer estimation performance than other methods when the covariate signal is weak. This suggests that when the covariate signal in the dataset is not strong compared to the spatial signal then it is not worth attempting any flexible covariate modeling. This is supported by our simulation studies. We hope our study provides inspiration for future spatial modeling studies with complex covariate patterns. Another potential extension is to consider non-Gaussian likelihood models where the outcome variable is discrete (i.e.,~counts and proportions) using the P\'{o}lyâ Gamma data augmentation technique \citep{polson2013bayesian}.

The key challenge that arises when one combines  spatial models with machine learning techniques, is attaching an appropriate measure of uncertainty. In general, uncertainty estimation is difficult with machine learning techniques, since the bootstrap does not work in many instances. For example, with sparse estimators this occurs because the limiting distribution is complex and may not be continuous, see \citealp{dezeure2015high}. When combining machine learning techniques with spatial models, this aspect becomes even more challenging. %Another challenge is the computation time -- BARTSIMP is not fast. 

\bibliographystyle{elsarticle-harv} 
\bibliography{main}

\newpage

\appendix




\end{document}

\endinput
%%
%% End of file `elsarticle-template-harv.tex'.


